\section{Introduction}

\looseness=-1
Voice has provided a convenient interface to early conversational systems, from Alexa\footnote{\url{https://www.alexa.com}} to Siri\footnote{\url{https://www.apple.com/siri}} and Google Assistant.\footnote{\url{https://assistant.google.com/}} In this context, a ``wake word'' spoken by the user typically triggers an automatic speech recognition (ASR) system which transcribes the subsequent user's request. Then, a natural language understanding (NLU) pipeline converts this query to a structured format used to produce a text answer through natural language generation (NLG). Eventually, a text-to-speech (TTS) system tells the answer back to the user. While this process can handle short, constrained interactions (e.g. triggering an action or retrieving a fact), the rise of large language models (LLMs)~\citep{gpt3,chinchilla,llama1} has called for a consequent extension of voice interfaces to multi-turn, open-ended conversations. A solution to this challenge is handling the NLU and NLG with an LLM, while the ASR and TTS provide the voice interface during the user's and the system's turn respectively~\citep{llama3herd}. This framework supports the current generation of spoken dialogue systems such as Gemini~\citep{team2023gemini} or ChatGPT.\footnote{\url{https://openai.com/index/chatgpt-can-now-see-hear-and-speak/}}%

\looseness=-1
Yet, the experience offered by these interfaces remains far from natural conversations. First, latency compounds along the many components of these pipelines, resulting in a typical global latency of several seconds. This is unlike natural conversations which demonstrate response times of a few hundred milliseconds. Second, as language understanding and generation happens in the textual domain, any non-written information is ignored by the model. This goes from paralinguistic information, such as emotion and accent, to non-speech audio, such as surrounding acoustic events. Finally, these models remain fundamentally turn-based, assuming that dialogue is a sequence of well-defined single-speaker segments. While this paradigm is suited to text dialogue, it falls short in modeling aspects of spoken conversations such as interruptions, overlapping speech--- which amounts for $10$ to $20\%$ of spoken time~\citep{overlap_stats_2} ---and backchanneling (i.e. non-interrupting interjections such as ``OK'' or ``I see'').

\looseness=-1
In this work we introduce \ours{}, a speech-text foundation model and real-time spoken dialogue system that aims at solving the aforementioned limitations: latency, textual information bottleneck and turn-based modeling. \ours{} augments a text LLM backbone with a smaller audio language model~\citep{audiolm,uniaudio} that ingests and predicts discrete audio units. This removes the information bottleneck of text by understanding inputs and generating outputs directly in the audio domain, while benefiting from the knowledge and reasoning abilities of the underlying text LLM. We extend previous work on audio language models and design a streaming, hierarchical architecture, with a theoretical latency of \SI{160}{\ms}---lower than the \SI{230}{\ms} average in natural conversations measured over 10 languages~\citep{Stivers2009UniversalsAC}. We furthermore introduce the first multi-stream audio language model, i.e. a model that explicitly processes the input and output audio streams jointly into two autoregressive token streams. This altogether removes the concept of speaker turn and thus allows training the model on natural conversations with arbitrary dynamics including overlap and interruptions. Our resulting model is the first \textit{full-duplex}--- it always listens and always generates sound, either speech or silence---real-time conversational LLM. We summarize our contributions below:
\begin{itemize}
\item We present \helium, a 7B-parameter text LLM that we pretrain on $2.1$T tokens of public English data. \hyperref[sec:helium]{Section \ref{sec:helium}} describes the architecture and training of the model, while \hyperref[sec:text_dataset]{Section \ref{sec:text_dataset}} provides details on the pretraining data collection and filtering.
\looseness=-1
\item We train \mimi, a neural audio codec~\citep{soundstream,encodec} that converts audio into the discrete tokens predicted by \ours{} and back, using residual vector quantization (RVQ). Audio language models typically combine such \textit{acoustic} tokens with \textit{semantic} tokens from a self-supervised speech model as it is necessary to produce intelligible speech in absence of text conditioning~\citep{audiolm}. We rather extend the approach of~\citet{zhang2024speechtokenizer} by distilling semantic information into the first level of acoustic tokens and introduce improved training tricks. \hyperref[sec:mimi]{Section~\ref{sec:mimi}} describes the architecture and training of \mimi while \hyperref[sec:mimi_eval]{Section~\ref{sec:mimi_eval}} details ablation studies.
\looseness=-1
\item We propose \ours{}, a new architecture for audio language modeling, which combines Helium with a smaller Transformer~\citep{attentionvaswani} model to predict audio tokens in a hierarchical and streaming fashion. We show how challenging it is for such unconditioned audio language models to generate intelligible speech, and we provide solutions that outperform the intelligibility and audio quality of non-streaming models while generating audio in a streaming fashion. We furthermore extend this architecture to model several audio streams in parallel, allowing for a conceptually and practically simple handling of full-duplex dialogues with arbitrary dynamics. \hyperref[sec:modeling]{Section \ref{sec:modeling}} describes this architecture.
\looseness=-1
\item In \hyperref[sec:interleaving]{Section~\ref{sec:interleaving}}, we introduce \interleaving, a new training and inference setup for audio language models that significantly improves the factuality and linguistic quality of generated speech by predicting time-aligned text tokens before audio tokens. \ours{} is a speech-to-speech model as it allows reasoning about non-linguistic information, both from the user audio and from \ours{}'s audio. Yet, this is not incompatible with \ours{} producing text along its speech output. Based on the past observation~\citep{audiolm,zhang2024speechtokenizer} that coarse-to-fine generation (from semantic to acoustic tokens) is critical to generating consistent speech, we extend this hierarchy to using text tokens as a per-timestep prefix to the semantic token. Our experiments show that not only this drastically improves the length and quality of generated speech, but we also show how forcing a delay between text and audio tokens allows deriving streaming ASR and streaming TTS from a \ours{} model.
\item We evaluate all components of \ours{} along several axes, including text understanding, speech intelligibility and consistency, audio quality and spoken question answering. Our experiments, reported in \hyperref[sec:evaluation]{Section~\ref{sec:evaluation}}, show that our model is state of the art among existing speech-text models for speech modeling and spoken question answering while being streaming compatible and able to model several minutes of context (\SI{5}{\minute} in our experiments).
\end{itemize}
We encourage the reader to talk to \ours{} using our web demo.\footnote{\url{https://moshi.chat/}}


\newpage
\section{Related Work}

\paragraph{Audio Language Modeling.}
\looseness=-1
Early developments in speech foundation models have improved speech understanding across many discriminative tasks, from automatic speech recognition (ASR)~\citep{wav2vec2,whisper,google_usm} to speaker verification~\citep{wavlm} and speech classification~\citep{superb}. A key factor in this development is self-supervised learning~\citep{hubert, wav2vec2,wavlm} which allows learning generic, discriminative speech representations. As these speech understanding models build on previous work done on masked language modeling for text~\citep{bert}, generative text pretraining~\citep{gpt} has similarly inspired a large family of speech generation models. In particular,~\citet{gslm} propose quantizing aforementioned self-supervised representations. The resulting discrete \textit{audio tokens} represent a speech segment as a sequence of categorical variables, thus casting speech generation as a language modeling task. AudioLM \citep{audiolm} furthermore combines these \textit{semantic} tokens with \textit{acoustic} tokens from a neural audio codec~\citep{soundstream}, which allows for modeling arbitrary voices, recording conditions and non-speech sounds.  These audio language models have redefined the state of the art in speech generation, from text-to-speech~\citep{valle, spear_tts} to speech-to-speech translation~\citep{audiopalm, gemini_1.5} and speech enhancement~\citep{uniaudio}. Beyond these supervised tasks, a parallel line of work has explored training and scaling unsupervised audio-only models, trained for autoregressive speech generation~\citep{zerospeech21,gslm,audiolm}. The abilities of these models have progressively expanded, from generating short sentences in a single speaker voice~\citep{gslm} to producing meaningful and consistent speech continuations across dozens of seconds in arbitrary voices and conditions~\citep{audiolm}, thanks to a hierarchical modeling of semantic and acoustic tokens. A main challenge is that audio requires the modeling of long sequences, up to a few minutes, to produce meaningful and exploitable outputs. However, latent representations for audio are typically less compact than equivalent representations for text. Thus, discrete representations from neural audio codecs require multiple predictions per timestep when modeled autoregressively. \citep{liu2023audioldm} and \citep{evans2024stable} use latent diffusion~\citep{ho2020denoising} for general audio and music modeling
to alleviate the need for hierarchical discrete tokens. However, these methods cannot be used in a streaming fashion, and it is unclear whether they could generate consistent speech. \citet{musicgen} instead show that the number of auto-regressive steps can be reduced by introducing a delay between the different levels of tokens, and performing parallel prediction over them. Inspired by the RQ-Transformer method by~\citet{rqtransformer} and the hierarchical MegaByte transformer model~\citep{yu2024megabyte},
\citet{uniaudio} and \citet{zhu2024generativepretrainedspeechlanguage} leverage a smaller nested transformer to model the different tokens at a single time step. In this work, we extend these previous works to push the limits of autoregressive speech generation by proposing a scalable hierarchical modeling of audio tokens which can handle several minutes of context while generating audio in real time. Still, while speech-only models learn linguistic structure---lexicon, syntax, semantics--- from raw speech~\citep{zerospeech21}, they typically demonstrate poor-to-nonexistent factual knowledge and reasoning abilities. This has led to the development of speech-text models, intended to combine the knowledge and reasoning abilities of text models with the generative power of audio models.


\paragraph{Speech-text Models.}
\looseness=-1
Such models typically start from a pretrained text language model and either finetune it to predict audio~\citep{twist}, or propose a speech-text finetuning task~\citep{audiopalm,voxtlm,spectron,spiritlm,mitsui2024pslmparallelgenerationtext,zhang2024speechgpt}: For instance, AudioPALM~\citep{audiopalm} starts from a pretrained PALM~\citep{palm} model, and extends its text vocabulary with semantic audio tokens. Then, the model is trained for a mixture of speech-text tasks, including TTS, ASR and speech-to-speech translation. VoxTLM~\citep{voxtlm} adopts a similar approach for TTS and ASR. While these models are trained in a supervised fashion with specific input and output sequences, Spirit-LM~\citep{spiritlm} uses temporal alignment between speech and its transcript to perform modality switch (from speech tokens to text tokens, or conversely) inside a sequence. This allows the model to learn consistent internal representations of language regardless of it being represented as text or speech, as measured through commonsense evaluation. Another approach, adopted by Spectron~\citep{spectron}, SpeechGPT~\citep{speechgpt_emnlp} and PSLM~\citep{mitsui2024pslmparallelgenerationtext}, combines speech and text in a hierarchical manner rather than as interchangeable representations. Similar to how AudioLM~\citep{audiolm} decomposes speech generation into predicting semantic tokens and \textit{then} acoustic tokens, Spectron and SpeechGPT use a ``Chain-of-Modality'' and first produce an utterance as text tokens, subsequently used as a prefix to generate speech. This allows guiding speech generation with the output of an underlying text LLM, however this is fundamentally incompatible with live interactions as the model needs to produce an entire answer as text before it starts speaking. PSLM alleviates this limitation by modeling text and speech tokens in parallel. In this work, we propose \interleaving as a main architectural and training component to combine aligned text and speech data. \interleaving decomposes speech into a chain of text, semantic and acoustic tokens, and predicts this structured sequence in a hierarchical manner. Unlike Spirit-LM, this allows representing all utterances both as text and speech, rather than switching between modalities; In addition, the integration of acoustic tokens into the same generative model enables generating arbitrary voices and conditions, rather than a single speaker. Besides, this hierarchical modeling described in \hyperref[sec:interleaving]{Section~\ref{sec:interleaving}} allows decomposing the generation task without increasing the sequence length of the Transformer~\citep{attentionvaswani} outputs, unlike Chain-of-Modality, while benefiting from producing text a prefix to audio tokens rather than in parallel like PSLM. Moreover, \interleaving decomposes speech on a per-frame basis, which means that each prediction step outputs a speech frame. This is unlike Spectron and SpeechGPT which require generating a complete sequence as text before generating audio tokens, and this makes \ours{} compatible with real-time generation. Moreover, we show in \hyperref[sec:interleaving]{Section~\ref{sec:interleaving}} how \interleaving, when combined with a delay between token types, allows deriving streaming TTS and ASR systems from \ours{}. Finally, while Spectron, SpeechGPT and PSLM model both user and system speech and text tokens into a single stream, which requires properly segmented turns, \ours{} benefits from a novel multi-stream architectures which removes the concept of speaker turns and allows for modeling any type of overlap, interruptions and interjections.

\paragraph{Spoken Dialogue Models.}
\looseness=-1
Spoken dialogue is one of the less explored tasks in speech generation, as it requires addressing several challenges: 1) The model should run in real time and allow for long conversations in full-duplex---the model always listens and can speak at any moment; 2) it should be speech-to-speech to handle paralinguistic communication; 3) it should display knowledge and reasoning abilities that make it amenable to helpful and enjoyable conversations. Spectron benefits from its underlying text LLM (as measured by spoken question answering), however it is not compatible with real-time generation due to Chain-of-Modality. PSLM proposes generating speech and text tokens in parallel to reduce this latency, however it reduces the quality of answers, and the model still relies on ASR, which removes paralinguistic information. More importantly, these models cannot handle full-duplex communication, where there is no boundary between speaker turns, as any side of the conversation can be active at any time. An attempt at modeling these dynamics has been proposed by \citet{fullduplexfsm}, with an ASR system running in parallel to a text generator that feeds into a streaming TTS. While this allows modeling more complex scenarios than previous approaches, it still relies on a cascaded pipeline and models both the user's and the system's speech into a single token stream, which is challenging in presence of significant overlap. The only previous full-duplex dialogue system is dGSLM~\citep{nguyen-etal-2023-generative}, which models user and system speech as separate audio token streams and proposes a Siamese architecture to process both streams jointly. While dGSLM is full-duplex, it remains a proof-of-concept: it does not run in an online fashion, it does not benefit from the knowledge of a text language model, and it does not model acoustic information as it only models semantic tokens. \ours{} addresses these limitations altogether: by modeling two streams of semantic and acoustic tokens hierarchically, \ours{} is full duplex and can exploit all the information from the user (linguistic and non-linguistic) while producing speech in real time. Thanks to text pretraining and \interleaving, \ours{} benefits from the knowledge of its \helium backbone. Finally, as the model produces acoustic tokens along with text and semantic tokens, it can generate an arbitrary range of emotions, voices and acoustic conditions. To the best of our knowledge, \ours{} is the first audio language model that successfully addresses the many aforementioned challenges of spoken dialogue.
